\uuid{qWTG}
\exo7id{1253}
\titre{exo7 1253}
\auteur{ridde}
\organisation{exo7}
\datecreate{1999-11-01}
\isIndication{false}
\isCorrection{false}
\chapitre{Développement limité}
\sousChapitre{Applications}
\cpge{mpsi-mp2i.limites-continuite.b, pcsi.limites-continuite.b, ptsi.limites-continuite.b, mpsi-mp2i.analyse-asymptotique.b, pcsi.analyse-asymptotique.b, ptsi.analyse-asymptotique.b, mpsi-mp2i.derivabilite.a, pcsi.derivabilite.a, ptsi.derivabilite.a, bcpst1.limites-continuite.c, bcpst1.developpements-limites.a, bcpst1.developpements-limites.b}
\module{Analyse}
\niveau{L1}
\difficulte{}

\contenu{
\texte{
Soit $f (x) = (\cos x)^{\frac 1x}$ pour $x \in ]-\frac{\pi}2, \frac{\pi}2[-
\left\{ 0\right\}$.
}
\begin{enumerate}
    \item \question{Montrer que $f$ est prolongeable par continuit\'e en $0$.}
    \item \question{D\'eterminer un DL de $f$ en $0$ \`a l'ordre $2$.}
    \item \question{Etudier la d\'erivabilit\'e du prolongement de $f$.}
\end{enumerate}
}
